\documentclass[11pt,letterpaper]{article}
\usepackage[T1]{fontenc}
\usepackage[utf8]{inputenc}
\usepackage[margin=0.88in,top=0.84in,bottom=0.83in,headheight=13pt,headsep=20pt,footskip=28pt]{geometry}
\usepackage{newpxtext}
\usepackage{amsmath,amsthm}
\usepackage{newpxmath}
% newpxtext selects TeX Gyre Heros (qhv), scaled to .94, for sans serif.
\usepackage{microtype}
\usepackage{xcolor}
\definecolor{Forest}{HTML}{30392C}
\definecolor{Olive}{HTML}{687144}
\definecolor{Muted}{HTML}{65685F}
\definecolor{Sage}{HTML}{CBD0BC}
\definecolor{Pale}{HTML}{F2F4EB}
\usepackage{mathtools}
\usepackage{booktabs,tabularx,longtable,array}
\usepackage{enumitem}
\usepackage{titlesec}
\usepackage{fancyhdr}
\usepackage[most]{tcolorbox}
\usepackage{needspace}
\usepackage{xurl}
\usepackage[colorlinks=true,linkcolor=Olive,citecolor=Olive,urlcolor=Olive,bookmarksnumbered=true,pdfencoding=auto,psdextra]{hyperref}
\hypersetup{pdftitle={Large cardinals at the inconsistency frontier: thin sections and simultaneous saturation},pdfauthor={Prepared for Vladimir Reshetnikov},pdfsubject={A research continuation with detailed proofs of definable thin-section obstructions},pdfkeywords={ultraexacting, exacting, finite modification, complete section, ordinal definability, I0, Icarus}}
\urlstyle{same}
\setlength{\parindent}{0pt}
\setlength{\parskip}{5.1pt plus 1pt minus .4pt}
\linespread{1.035}
\setlength{\emergencystretch}{2em}
\setcounter{tocdepth}{1}
\setcounter{secnumdepth}{2}
\titleformat{\section}{\Large\bfseries\color{Forest}}{\thesection}{.6em}{}
\titleformat{\subsection}{\large\bfseries\color{Forest}}{\thesubsection}{.55em}{}
\titleformat{\subsubsection}{\normalsize\bfseries\color{Forest}}{}{0pt}{}
\titlespacing*{\section}{0pt}{18pt plus 3pt minus 2pt}{8pt}
\titlespacing*{\subsection}{0pt}{12pt plus 2pt minus 1pt}{5pt}
\titlespacing*{\subsubsection}{0pt}{9pt}{3pt}
\setlist[itemize]{leftmargin=1.4em,itemsep=2pt,topsep=3pt,parsep=0pt}
\setlist[enumerate]{leftmargin=1.7em,itemsep=3pt,topsep=4pt,parsep=0pt}
\pagestyle{fancy}
\fancyhf{}
\fancyhead[L]{\sffamily\fontsize{9}{11}\selectfont\color{Muted}LARGE CARDINALS AT THE INCONSISTENCY FRONTIER}
\fancyhead[R]{\sffamily\fontsize{9}{11}\selectfont\color{Muted}18 SEP 2026}
\fancyfoot[L]{\small\color{Muted}Thin sections and simultaneous saturation}
\fancyfoot[R]{\small\color{Muted}\thepage}
\renewcommand{\headrulewidth}{.35pt}
\renewcommand{\footrulewidth}{0pt}
\fancypagestyle{plain}{\fancyhf{}\fancyfoot[L]{\small\color{Muted}Thin sections and simultaneous saturation}\fancyfoot[R]{\small\color{Muted}\thepage}\renewcommand{\headrulewidth}{0pt}}
\tcbset{colback=Pale,colframe=Sage,colbacktitle=Sage,coltitle=Forest,fonttitle=\bfseries,boxrule=.5pt,arc=3pt,left=9pt,right=9pt,top=6pt,bottom=6pt,toptitle=3pt,bottomtitle=3pt,before skip=8pt,after skip=8pt,breakable}
\newtcolorbox{assessment}[1]{title={#1}}
\theoremstyle{definition}
\newtheorem{definition}{Definition}[section]
\newtheorem{theorem}[definition]{Theorem}
\newtheorem{proposition}[definition]{Proposition}
\newtheorem{lemma}[definition]{Lemma}
\newtheorem{corollary}[definition]{Corollary}
\newcommand{\ZFC}{\mathsf{ZFC}}
\newcommand{\ZF}{\mathsf{ZF}}
\newcommand{\AC}{\mathsf{AC}}
\newcommand{\DC}{\mathsf{DC}}
\newcommand{\HOD}{\mathrm{HOD}}
\newcommand{\OD}{\mathrm{OD}}
\newcommand{\HH}{\mathsf{HH}}
\newcommand{\Ex}{\operatorname{Ex}}
\newcommand{\UEx}{\operatorname{UEx}}
\newcommand{\Ext}{\operatorname{Ext}}
\newcommand{\SC}{\operatorname{SC}}
\newcommand{\CEx}{\operatorname{CEx}}
\newcommand{\Con}{\operatorname{Con}}
\newcommand{\crit}{\operatorname{crit}}
\newcommand{\cf}{\operatorname{cf}}
\newcommand{\ot}{\operatorname{ot}}
\newcommand{\Ord}{\mathrm{Ord}}
\newcommand{\Card}{\mathrm{Card}}
\newcommand{\rank}{\operatorname{rank}}
\newcommand{\id}{\mathrm{id}}
\newcommand{\restr}{\mathbin{\upharpoonright}}
\newcommand{\Izero}{\mathrm{I0}}
\newcommand{\Ione}{\mathrm{I1}}
\newcommand{\Itwo}{\mathrm{I2}}
\newcommand{\Ithree}{\mathrm{I3}}
\newcommand{\Isharp}{\mathrm{I0}^{\#}}
\newcommand{\Status}[1]{\textbf{Status.} #1}
\newcommand{\Priority}[1]{\textbf{Research priority.} #1}
\newcolumntype{Y}{>{\raggedright\arraybackslash}X}
\newcolumntype{P}[1]{>{\raggedright\arraybackslash}p{#1}}
\newcommand{\arxiv}[1]{\href{https://arxiv.org/abs/#1}{arXiv:#1}}
\newcommand{\doi}[1]{\href{https://doi.org/#1}{doi:\nolinkurl{#1}}}

\newtheorem{remark}[definition]{Remark}
\newtheorem{question}[definition]{Question}
\newtheorem{inputthm}{Input}
\renewcommand{\theinputthm}{\Alph{inputthm}}
\numberwithin{equation}{section}
\newcommand{\Pow}{\mathcal P}
\newcommand{\ran}{\operatorname{ran}}
\newcommand{\tc}{\operatorname{tc}}
\newcommand{\Good}{\mathcal G}
\newcommand{\Thin}{\operatorname{Thin}}
\newcommand{\cs}{\mathrm{CS}}
\newcommand{\eqfin}{\mathrel{=^{*}}}
\newcommand{\Tbase}{T_{\partial}}
\newcommand{\Tsat}{T_{\mathrm{sat}}}
\newcommand{\defcode}{\operatorname{code}}
\newcommand{\Sat}{\operatorname{Sat}}
\begin{document}
\thispagestyle{empty}
{\sffamily\bfseries\small\color{Olive}SET THEORY\quad /\quad RESEARCH CONTINUATION}

\vspace{13pt}
{\fontsize{28}{31}\selectfont\bfseries\color{Forest}Large cardinals at the\\ inconsistency frontier:\\ thin sections and saturation\par}
\vspace{10pt}
{\fontsize{14}{18}\selectfont A cardinal-sized strengthening of the finite-transversal obstruction,\\ with detailed proofs and an $\Izero$-equiconsistency refinement\par}
\vspace{14pt}
{\color{Muted}Prepared for Vladimir Reshetnikov\hfill 18 September 2026\par}
\vspace{3pt}
{\color{Sage}\hrule height .6pt}
\vspace{12pt}

\textbf{Abstract.}
Let $D_\lambda$ be the cofinal subsets of $\lambda$ of order type $\omega$, and let
$Q_\lambda=D_\lambda/{\eqfin}$, where $\eqfin$ is equality modulo finite difference.
We strengthen the finite-valued transversal obstruction to arbitrary fibers of
cardinality below $\lambda$, with a common family of classes for all low-rank
parameters. If $\lambda$ is ultraexacting, then at least $\lambda$ classes
$q\in Q_\lambda$ satisfy
\[
  T\in\OD_{V_\lambda},\quad T\subseteq D_\lambda
  \quad\Longrightarrow\quad |T\cap q|\in\{0,\lambda\}.
\]
The same classes work for every such $T$, even with different parameters from
$V_\lambda$ and arbitrarily large ordinal parameters. The proof combines fixed
small cofinal unions, a fixed bijection $\lambda\to V_\lambda$, and uniform
bounded-rank correctness of ordinal definability. Both the fiber-cardinality
cutoff and the uniform parameter-rank cutoff are sharp. We also exclude
nonempty $\OD_{V_\lambda}$ families of fewer than $\lambda$ thin complete
sections, without requiring their members definable or uniformly bounding
their fibers. Further results give a preserved-predicate obstruction to an
Icarus embedding, bounded-rank separation under ordinary exactingness, and
an $\Izero$-equiconsistent model with this saturation, $V_\lambda\subseteq\HOD$,
and definable choice on the bounded quotient.

\begin{assessment}{What is established, and what is not}
The main conclusions are \textbf{conditional inconsistency theorems in $\ZFC$}:
ultraexactingness is incompatible with the stated definable thinness principles.
Together with results of Aguilera--Bagaria--Goldberg--L\"ucke, they give
an $\Izero$-equiconsistent model with the new saturation properties and
$V_\lambda\subseteq\HOD$.

The new statements strengthen the supplied reports; worldwide priority has not
been established. The proofs are conventional and have not been independently
refereed or machine-checked. No unconditional inconsistency of a standard
large-cardinal hypothesis is claimed.
\end{assessment}

\newpage
\tableofcontents

\newpage
\section{The result and its relationship to the supplied research}\label{sec:overview}

The two supplied documents play different roles. The \emph{Unified Report}
\cite{unified} maps several consistency boundaries. The \emph{Synthesis}
\cite{synthesis} consolidates nine continuations, including cofinality barriers
for cover exactingness, finite-choice obstructions at ultraexacting cardinals,
and consistency calibrations. Repeating those results under new names would not
advance the investigation. This continuation concentrates on the finite-valued
transversal theorem in its Section~9 and asks how much of its finiteness
hypothesis is genuinely needed.

The essential observation is elementary but stronger than finite periodicity:
\emph{the union of fewer than $\lambda$ countable sets has cardinality below
$\lambda$}. If such a family of cofinal sets is preserved by an exacting
embedding, its union is a preserved short cofinal subset of $\lambda$,
contradicting the critical-sequence obstruction. This works even though
$\cf(\lambda)=\omega$; no regularity of $\lambda$ in $V$ is being assumed.

The harder part is the quantifier order. A particular ordinal-definable set
need not be fixed by a chosen embedding, since the ordinal parameters in its
definition need not be fixed. We handle all definitions simultaneously by
stabilizing bounded-rank ordinal definability before choosing the embedding,
and then selecting a canonical least counterexample from the fixed parameters.
Section~\ref{sec:od} supplies the full bookkeeping.

\subsection{Objects and terminology}

For an uncountable cardinal $\lambda$ with $\cf(\lambda)=\omega$, set
\begin{align}
 D_\lambda&=\{a\subseteq\lambda:\ot(a)=\omega\text{ and }\sup a=\lambda\},\\
 a\eqfin b&\quad\Longleftrightarrow\quad |a\triangle b|<\omega,\\
 [a]&=\{b\in D_\lambda:b\eqfin a\},\qquad
 Q_\lambda=\{[a]:a\in D_\lambda\}.
\end{align}
Here $\OD_p$ means definability from the \emph{single set parameter} $p$ and
finitely many ordinals. By contrast, $\OD_{V_\lambda}$ allows finitely many
parameters that are elements of $V_\lambda$. A finite tuple of such parameters
can be coded by one element of $V_\lambda$.

\begin{definition}\label{def:thin}
A \emph{complete section} is a set $T\subseteq D_\lambda$ meeting every class in
$Q_\lambda$. It is \emph{thin} if
\[
  0<|T\cap q|<\lambda\qquad(q\in Q_\lambda).
\]
There is no requirement that one cardinal below $\lambda$ bound all these
intersections. A transversal is a complete section meeting each class exactly
once.
\end{definition}

\begin{definition}\label{def:good}
For any set parameter $p$, define the set of \emph{fully saturated classes}
\[
 \Good_p(\lambda)=
 \left\{q\in Q_\lambda:
   \forall T\in\OD_p\,
   \bigl(T\subseteq D_\lambda\ \Longrightarrow\ |T\cap q|\in\{0,\lambda\}\bigr)
 \right\}.
\]
``Saturated'' in this report refers only to this explicitly defined property,
not to saturation of an ideal, a model, or a Boolean algebra.
\end{definition}

\begin{definition}[Simultaneity over all low-rank parameters]\label{def:good-low}
Let
\[
 \Good_{\mathrm{low}}(\lambda)=
 \left\{q\in Q_\lambda:
   \forall T\in\OD_{V_\lambda}\,
   \bigl(T\subseteq D_\lambda\ \Longrightarrow\ |T\cap q|\in\{0,\lambda\}\bigr)
 \right\}.
\]
Thus $\Good_{\mathrm{low}}(\lambda)=\bigcap_{p\in V_\lambda}\Good_p(\lambda)$.
\end{definition}

\begin{theorem}[Main saturation theorem]\label{thm:main-intro}
In $\ZFC$, if $\lambda$ is ultraexacting, then
\[
                       |\Good_{\mathrm{low}}(\lambda)|\ge\lambda.
\]
In particular, no thin complete section belongs to $\OD_{V_\lambda}$.
\end{theorem}

The theorem will be proved as Theorem~\ref{thm:main}. Its first conclusion
supplies a common set of classes for \emph{all} definitions from elements of
$V_\lambda$, not just all definitions with one fixed parameter. The proof
constructs the common classes directly: it does not infer a large intersection
merely from knowing that each $\Good_p(\lambda)$ is large.

\subsection{An explicit conditional inconsistency statement}

Write $\mathsf{ThinCS}(\lambda)$ for
\[
 \exists p\in V_\lambda\ \exists T\in\OD_p\quad
    \bigl(T\subseteq D_\lambda\ \wedge\
          \forall q\in Q_\lambda\ (0<|T\cap q|<\lambda)\bigr).
\]
Then the central negative result is the ordinary first-order implication
\begin{equation}\label{eq:inconsistency}
 \ZFC\ \vdash\
 \neg\exists\lambda\,
       \bigl(\UEx(\lambda)\wedge\mathsf{ThinCS}(\lambda)\bigr).
\end{equation}
Ordinal definability is first-order definable, so this is not an unspecified
scheme with a new truth predicate. An even stronger refutation, involving
small families of potentially nondefinable sections, appears in
Section~\ref{sec:families}.

\section{Framework and external inputs}\label{sec:framework}

We work in an ambient universe satisfying $\ZFC$. Cardinalities, power sets,
and cofinality are ambient unless an inner model is explicitly indicated.
All witnesses in the main argument are \emph{set} embeddings between set
structures; their elementarity is expressed using the usual satisfaction
relation for sets. Their domains need not be transitive.

\begin{definition}[Exacting and ultraexacting]\label{def:uex}
A cardinal $\lambda$ is exacting if, at every sufficiently large rank height
$\zeta$, there are
\[
 X\prec V_\zeta,\qquad V_\lambda\cup\{\lambda\}\subseteq X,
 \qquad j:X\longrightarrow V_\zeta
\]
with $j$ elementary, $j(\lambda)=\lambda$, and $j\restr\lambda\ne\id$.
It is ultraexacting if these witnesses can also satisfy
$j\restr V_\lambda\in X$. We use the equivalent all-height/high-critical-point
formulation in Input~\ref{in:high}.
\end{definition}

The primary literature provides the following inputs. They are separated
from the new proofs so that the external dependencies are visible.

\begin{inputthm}[Local Kunen inconsistency]\label{in:kunen}
There is no nontrivial elementary self-embedding of $V_{\mu+2}$ in $\ZFC$.
\cite{kunen,cover}
\end{inputthm}

\begin{inputthm}[Choice of witnesses]\label{in:high}
If $\lambda$ is ultraexacting, then for any sufficiently large height $\zeta$
and any $\beta<\lambda$, an ultraexacting witness can be chosen with
$\crit(j)>\beta$. The corresponding high-critical-point formulation also
holds for ordinary exacting witnesses.
\cite[Lemmas 2.3 and 3.2]{abl}
\end{inputthm}

The strict inequality is obtained from the cited identity-on-an-initial-segment
form by applying it to $\beta+1$. In particular, choosing the critical point
above the rank of $p\in V_\lambda$ makes $j(p)=p$.

\begin{inputthm}[Consistency comparison and coding]\label{in:cons}
The theories $\ZFC+\exists\lambda\,\UEx(\lambda)$ and $\ZFC+\Izero$ are
equiconsistent. Moreover, an ultraexacting $\lambda$ can be preserved by a set
forcing that makes $V_\lambda\subseteq\HOD$.
\cite[Theorem 4.5 and Proposition 3.9]{abgl}
\end{inputthm}

These sources were consulted directly for the stated theorem and lemma
numbers. The argument below does not use the synthesis's
$\mathrm{I3}_{\mathrm{wf}}(0)$ forcing construction, its high-completeness
core theorem, or its unbounded-range obstruction for class embeddings.
Those separate lines are not dependencies of the results below.

\subsection{What ``detailed proofs'' covers here}

Every new implication is proved below, including the canonical-definition
argument, the construction of many fixed equivalence classes, the singular
cardinal estimate for nonuniform families, and both directions of the
consistency refinement. Input~\ref{in:cons} is an imported large-cardinal
comparison, not reproved or presented as a new comparison. Ordinary background
facts about set satisfaction, rank, elementary embeddings, and finite coding
are used explicitly where needed.

\section{The local obstruction: fixed small cofinal unions}\label{sec:local}

\begin{lemma}[Critical-sequence normalization]\label{lem:critical}
Let $j:X\to V_\zeta$ be an exacting witness at $\lambda$, where
$\zeta>\lambda+\omega$. Then $e=j\restr V_\lambda$ is a nontrivial elementary
self-embedding of $V_\lambda$. If
\[
 \kappa=\crit(e),\qquad \kappa_n=e^n(\kappa)\quad(n<\omega),
\]
then
\[
 \sup_{n<\omega}\kappa_n=\lambda,
 \qquad e(\xi)>\xi\quad(\kappa\le\xi<\lambda).
\]
Also $j$ fixes every element of $V_\kappa$.
\end{lemma}

\begin{proof}
The rank $V_\lambda$ is definable from $\lambda$. Since $\lambda\in X$,
$V_\lambda\in X$; since $V_\lambda\subseteq X$ as well, satisfaction in that
rank is computed over the full structure, not a truncated domain. Relativizing
any formula to $V_\lambda$ and using $j(\lambda)=\lambda$ therefore proves the
elementarity of $e$.

Its critical point is a cardinal, and is measurable by the usual derived
ultrafilter
\[
 U=\{A\subseteq\kappa:\kappa\in e(A)\}.
\]
All subsets and the short sequences needed to verify $\kappa$-completeness have
rank below $\lambda$. Thus $\kappa$ is strongly inaccessible. Elementarity gives
that each $\kappa_n$ is strongly inaccessible as well. In particular the
sequence is strictly increasing and bounded by $\lambda$.

Suppose $\mu=\sup_n\kappa_n<\lambda$. The whole sequence then has rank below
$\lambda$; applying $e$ shifts it by one place, so $e(\mu)=\mu$. Since $\lambda$
is a cardinal and $\mu<\lambda$, the rank $V_{\mu+2}$ is still below $\lambda$.
The restriction of $e$ to this definable, fixed rank is a nontrivial elementary
self-embedding, contrary to Input~\ref{in:kunen}. Hence the supremum is
$\lambda$.

For $\kappa\le\xi<\lambda$, choose $n$ with
$\kappa_n\le\xi<\kappa_{n+1}$. Then
$e(\xi)\ge e(\kappa_n)=\kappa_{n+1}>\xi$.
Finally, the usual rank induction below the inaccessible $\kappa$ proves that
$e\restr V_\kappa=\id$: each set of rank below $\kappa$ has size below
$\kappa$, so its image is its pointwise image; apply the induction hypothesis
to its elements. This also gives the last assertion for $j$.
\end{proof}

In particular, $\cf^V(\lambda)=\omega$ and $\lambda$ is a strong limit, as a
limit of inaccessible cardinals. The only fixed ordinals below $\lambda$ are
those below $\kappa$.

\begin{lemma}[Fixed small sets of ordinals]\label{lem:fixed-ordinals}
In the setting of Lemma~\ref{lem:critical}, suppose $A\in X$,
$A\subseteq\lambda$, $j(A)=A$, and $|A|<\lambda$. Then $A\subseteq\kappa$.
Consequently, no fixed cofinal subset of $\lambda$ has cardinality below
$\lambda$.
\end{lemma}

\begin{proof}
Put $\tau=\ot(A)$. Since $A$ is a subset of the cardinal $\lambda$,
$\tau\le\lambda$; the size assumption excludes equality, so $\tau<\lambda$.
The order type and the unique increasing enumeration $h:\tau\to A$ are
definable from $A$, and therefore belong to $X$. Because $A$ is fixed,
$j(\tau)=\tau$ and $j(h)=h$. Lemma~\ref{lem:critical} gives $\tau<\kappa$.
For every $\xi<\tau$, all the following evaluations are defined:
\[
 j(h(\xi))=j(h)(j(\xi))=h(\xi).
\]
Here $\xi\in X$ because $\xi<\lambda$, and $h(\xi)\in X$ by elementary
closure under function evaluation. Thus every member of $A$ is a fixed
ordinal below $\lambda$, and hence is below $\kappa$.
\end{proof}

\begin{lemma}[Fixed small families of cofinal sets]\label{lem:small-family}
In the same setting, there is no $B\in X$ satisfying
\[
 j(B)=B,\qquad \varnothing\ne B\subseteq D_\lambda,
 \qquad |B|<\lambda.
\]
\end{lemma}

\begin{proof}
The set $A=\bigcup B$ belongs to $X$ and is fixed by $j$, since union is a
uniquely definable operation. It is cofinal in $\lambda$, because any member
of $B$ is cofinal. In the ambient universe, Choice gives
\[
 |A|\le |B|\cdot\aleph_0\le\max\{|B|,\aleph_0\}<\lambda.
\]
This contradicts Lemma~\ref{lem:fixed-ordinals}.
\end{proof}

\begin{remark}[Why singularity causes no problem here]\label{rem:singular}
We are taking a union of a fixed number $\mu<\lambda$ of \emph{countable}
sets, not a union of $\mu$ sets whose cardinalities approach $\lambda$.
The estimate $\mu\cdot\aleph_0<\lambda$ is valid for every uncountable
cardinal $\lambda$. Section~\ref{sec:families} treats the genuinely nonuniform
case separately rather than silently assuming regularity.
\end{remark}

\section{Ordinal parameters and uniform bounded-rank correctness}\label{sec:od}

We need one precise device for turning a possibly non-fixed definable object
into a fixed least counterexample. It is useful to make the device uniform
before the critical-sequence class is known.

\subsection{A canonical order on $\OD_p$}

Fix a set $p$. A local definition code consists of an ordinal $\theta$ with
$p\in V_\theta$, a formula number, and a finite tuple of ordinals below
$\theta$. The code defines $x$ when $x$ is the unique element of $V_\theta$
satisfying that formula over $(V_\theta,\in)$ with those parameters and $p$.
This is a first-order relation in the ambient universe, since the satisfaction
relation is for a \emph{set} structure.

Order these codes first by $\theta$, and, at a fixed $\theta$, by a fixed
wellordering of formula numbers, finite lengths, and ordinal tuples. Each
stage is a set. Every $x\in\OD_p$ has such a code: apply reflection to a global
definition of $x$ and its uniqueness. Conversely, every locally coded set is
in $\OD_p$, since $\theta$ is an allowed ordinal parameter. Ordering sets by
their least codes gives a definable, set-like class wellorder $\prec_p$ of
$\OD_p$.

\begin{lemma}[Uniform bounded-rank stabilization]\label{lem:stabilization}
Fix a set $P$ of possible parameters and an ordinal $\delta$. There is a
height $\eta$ such that, at every sufficiently large limit cardinal
$\zeta>\eta$, simultaneously for every $p\in P$, $V_\zeta$ computes correctly
both
\[
             \OD_p\cap V_\delta
 \quad\text{and}\quad
             \prec_p\restr(\OD_p\cap V_\delta).
\]
In particular, one can take $P=\{p\}$ for a single prescribed parameter, or
$P=V_{\lambda+1}$ before the embedding is chosen.
\end{lemma}

\begin{proof}
The code relation defines $\OD_p$ uniformly in $p$. Thus
\[
             S=\{(p,x)\in P\times V_\delta:x\in\OD_p\}
\]
is a set by Separation. To each $(p,x)\in S$ assign the unique least code
of $x$ relative to $p$. Replacement gives a set of all these codes and a
bound on their heights and ranks. Choose $\eta$ above this bound, above the
ranks of $P$ and $V_\delta$, and above the finite coding overheads.

For a large enough limit cardinal $\zeta>\eta$, all the least codes under
consideration are present in $V_\zeta$, as are the rank structures and set
satisfaction relations used in those codes. Evaluation of a local code is
absolute. Hence every actual $\OD_p$ member of $V_\delta$ is recognized,
uniformly for $p\in P$. Conversely, a code used by $V_\zeta$ is a genuine
code in $V$: its structure is the full $V_\theta$, and finite-formula
satisfaction is absolute. There are no false positives. All codes preceding
a least code have no larger height, so no earlier competing code is omitted.
The least codes and their order are therefore correct, simultaneously for
all the indicated parameters.
\end{proof}

The lemma is not a truth predicate for $V$. It bounds a \emph{set} of local
codes for a previously fixed rank and a \emph{set} of possible parameters.
Its height can be arbitrarily large; ultraexactingness supplies witnesses
at the height chosen afterwards. The uniformity in $p$ will be essential:
the parameter coding $V_\lambda$ depends on the subsequently chosen embedding.

\begin{lemma}[Fixed canonical counterexamples]\label{lem:canonical}
Suppose $p,\lambda,q\in X$, $j(p)=p$, $j(\lambda)=\lambda$, and $j(q)=q$.
Let $\Phi(T,p,\lambda,q)$ describe objects $T$ in a fixed bounded rank, and
assume $V_\zeta$ is correct about $\Phi$, $\OD_p$, and $\prec_p$ on those
objects. If some $T\in\OD_p$ satisfies $\Phi$, then the $\prec_p$-least such
$T$ belongs to $X$ and is fixed by $j$.
\end{lemma}

\begin{proof}
Correctness makes the actual least object the unique solution of the same
formula in $V_\zeta$. By $X\prec V_\zeta$, that unique solution belongs to
$X$. Apply $j$ to its defining property. All displayed parameters are fixed,
so $j(T)$ is the same unique solution in $V_\zeta$. Hence $j(T)=T$.
\end{proof}

For our applications, $\Phi$ is built from $D_\lambda$, $Q_\lambda$, finite
symmetric difference, and cardinality below $\lambda$. At a height above
$\lambda+\omega$, all relevant sets, increasing enumerations, and possible
bijections with ordinals below $\lambda$ are present. These properties are
therefore absolute. Lemma~\ref{lem:stabilization}, with $\delta=\lambda+12$ and a suitable
set $P$, supplies the remaining correctness uniformly for all candidate
classes $q$ and all allowed coding parameters that the subsequently chosen
embedding might produce.

\needspace{10\baselineskip}
\section{Many fixed classes and simultaneous saturation}\label{sec:saturation}

\begin{lemma}[The size of a finite-modification class]\label{lem:class-size}
Every $q\in Q_\lambda$ has cardinality exactly $\lambda$.
\end{lemma}

\begin{proof}
For $q=[a]$, every member is determined by removing finitely many elements
of $a$ and adjoining finitely many ordinals below $\lambda$. There are at
most $\lambda$ such choices. Conversely, the sets $a\cup\{\xi\}$, for
$\xi\in\lambda\setminus a$, are distinct members of $[a]$. Adjoining finitely
many ordinals to an increasing cofinal $\omega$-set leaves its order type
$\omega$: below any fixed bound only finitely many elements occur. There are
$\lambda$ choices of $\xi$, so the lower bound follows.
\end{proof}

\begin{lemma}[Internal shifted critical sequences]\label{lem:shifted}
Let $j:X\to V_\zeta$ be an ultraexacting witness as above, and put
$\kappa=\crit(j)$ and $\kappa_n=j^n(\kappa)$. For every $t<\kappa$, the set
\[
       a_t=\{\kappa_n+t:n<\omega\}
\]
and its class $q_t=[a_t]$ belong to $X$. Moreover,
\[
 \begin{gathered}
 j(a_t)=a_t\setminus\{\kappa+t\},\qquad j(q_t)=q_t,\\
 a_t\cap a_s=\varnothing\quad(t\ne s<\kappa).
\end{gathered}
\]
\end{lemma}

\begin{proof}
The restriction $e=j\restr V_\lambda$ is an element of $X$. Its critical
point and its finite iterates on that critical point are uniformly definable
from $e$. The recursion of length $\omega$ is available in $V_\zeta$, and its
unique resulting sequence therefore belongs to $X$. It is the actual critical
sequence: all its values and finite evaluations occur in the full rank
$V_\lambda\subseteq X$.

Since $t<\kappa$, it belongs to $X$ and is fixed. Consequently $a_t$ and its
increasing enumeration belong to $X$. For a countable set with an enumeration
in $X$, elementarity and the fact that $j$ fixes every natural number imply
that its image is its pointwise image. Using preservation of ordinal addition,
\[
 j(a_t)=\{j(\kappa_n+t):n<\omega\}
       =\{\kappa_{n+1}+t:n<\omega\}.
\]
This is the stated tail. Since the equivalence class is definable from
$a_t$ and the fixed $\lambda$, it is in $X$, and
$j(q_t)=[j(a_t)]=[a_t]$.

Finally, the $\kappa_n$ are increasing infinite cardinals. For $t<\kappa$,
$\kappa_n+t<\kappa_n^+\le\kappa_{n+1}$. Thus the different $n$ lie in disjoint
intervals, and addition on the right is strictly increasing in $t$ within an
interval. The sets $a_t$ are pairwise disjoint and cofinal, so their classes
are pairwise distinct.
\end{proof}

\begin{lemma}[An internally available fixed code for the whole lower rank]\label{lem:fixed-code}
For every ultraexacting witness as in Lemma~\ref{lem:shifted}, there is a
bijection
\[
                    b:\lambda\longrightarrow V_\lambda
\]
such that $b\in X\cap V_{\lambda+1}$ and $j(b)=b$. Consequently,
\[
                        \OD_{V_\lambda}\subseteq\OD_b.
\]
This fixed-code construction is the same underlying device as the preserved
wellordering used in \cite[Proposition 3.9]{abgl}; its proof is given here.
\end{lemma}

\begin{proof}
Let $\kappa=\crit(j)$. Since $\kappa$ is inaccessible,
$|V_\kappa|=\kappa$. Choose a bijection $b_0:\kappa\to V_\kappa$.
Its graph has rank below $\lambda$, so $b_0\in V_\lambda\subseteq X$.
Put $e=j\restr V_\lambda$ and $b_n=e^n(b_0)$. Each $b_n$ is a bijection
$\kappa_n\to V_{\kappa_n}$.

For $\xi<\kappa$, both $\xi$ and $b_0(\xi)$ are fixed by $j$, so
\[
             b_1(\xi)=j(b_0)(\xi)=j(b_0(\xi))=b_0(\xi).
\]
Thus $b_0\subseteq b_1$. Applying $e$ successively shows
$b_n\subseteq b_{n+1}$ for every $n$. Since $e,b_0\in X$, the sequence
$\langle b_n:n<\omega\rangle$ is uniquely definable by recursion and belongs
to $X$. Its union $b$ belongs to $X$. The domains increase to $\lambda$ and
the ranges increase to $V_\lambda$, so $b$ is a bijection between them.

The countable sequence is in $X$, its indices are fixed, and
$j(b_n)=b_{n+1}$. Therefore
\[
                  j(b)=\bigcup_{n<\omega}b_{n+1}
                      =\bigcup_{n<\omega}b_n=b.
\]
Each ordered pair in the graph of $b$ has rank below the limit cardinal
$\lambda$, whence $b\subseteq V_\lambda$ and $b\in V_{\lambda+1}$.

Finally, every $p\in V_\lambda$ equals $b(\xi)$ for some ordinal
$\xi<\lambda$. A definition using finitely many such parameters can replace
each one by its uniquely specified value under $b$. The result is a
definition from $b$ and finitely many ordinals. This proves the inclusion.
Notice that the coding parameter $b$ need not itself belong to $V_\lambda$;
its fixedness is proved directly, not inferred from its rank.
\end{proof}

\begin{theorem}[Simultaneous full-sized traces for all low-rank parameters]\label{thm:main}
If $\lambda$ is ultraexacting, then
\[
                         |\Good_{\mathrm{low}}(\lambda)|\ge\lambda.
\]
Thus at least $\lambda$ classes have only empty or full-sized intersections
with every $\OD_{V_\lambda}$ subset of $D_\lambda$.
\end{theorem}

\begin{proof}
Fix an arbitrary $\beta<\lambda$. Apply
Lemma~\ref{lem:stabilization} with
\[
                       P=V_{\lambda+1},\qquad \delta=\lambda+12,
\]
and choose a sufficiently large limit cardinal $\zeta$ above its bound.
Use Input~\ref{in:high} to obtain an ultraexacting witness
$j:X\to V_\zeta$ with $\kappa=\crit(j)>\beta$.
Lemma~\ref{lem:fixed-code} supplies a fixed bijection
$b\in X\cap V_{\lambda+1}$ coding all of $V_\lambda$.
For each $t<\kappa$, Lemma~\ref{lem:shifted} supplies a fixed class
$q_t\in X$.

We claim that $q_t\in\Good_b(\lambda)$ for each $t<\kappa$. If not, the
fact that $|q_t|=\lambda$ gives some $T\in\OD_b$ satisfying
\begin{equation}\label{eq:bad-trace}
                  T\subseteq D_\lambda,
                  \qquad 0<|T\cap q_t|<\lambda.
\end{equation}
Although $b$ was unknown when $\zeta$ was chosen, it belongs to the set $P$
of parameters for which correctness was made uniform. Thus
$V_\zeta$ is correct about $\OD_b$, its canonical order, and
\eqref{eq:bad-trace}. Since $b,\lambda,q_t$ are fixed, the canonical least
witness $T_*$ is a fixed member of $X$, by
Lemma~\ref{lem:canonical}. Then
\[
                           B=T_*\cap q_t
\]
is a fixed member of $X$, is a nonempty family of members of $D_\lambda$,
and has cardinality below $\lambda$. This contradicts
Lemma~\ref{lem:small-family}.

Consequently all the $q_t$ are in $\Good_b(\lambda)$. The inclusion
$\OD_{V_\lambda}\subseteq\OD_b$ implies
\[
               \Good_b(\lambda)\subseteq\Good_{\mathrm{low}}(\lambda).
\]
The $\kappa$ classes $q_t$ are pairwise distinct, so
$|\Good_{\mathrm{low}}(\lambda)|\ge\kappa>\beta$. Since $\beta<\lambda$
was arbitrary, the desired bound follows. Explicitly, a hypothetical
cardinality $\mu<\lambda$ would be contradicted by the same construction
with $\beta=\mu$.
\end{proof}

\begin{remark}[How the quantifiers are handled]
No intersection principle is being assumed. A single fixed coding parameter
$b$ covers \emph{every} low-rank parameter. Correctness is uniform over all
possible $b$ before the embedding is chosen. The least-counterexample
argument then works for arbitrary ordinal parameters, even ones that $j$
moves. These three ingredients give the common set of classes directly.
\end{remark}

\begin{corollary}[Fixed parameters and bounded ranks]\label{cor:bounded-params}
For each $p\in V_\lambda$, $|\Good_p(\lambda)|\ge\lambda$. For every
$\alpha<\lambda$, at least $\lambda$ classes also have only empty or
$\lambda$-sized intersections with all $\OD_{V_\alpha}$ subsets of
$D_\lambda$. In both assertions the same set
$\Good_{\mathrm{low}}(\lambda)$ works simultaneously.
\end{corollary}

\begin{proof}
Both $\OD_p$ and $\OD_{V_\alpha}$ are subclasses of
$\OD_{V_\lambda}$. Apply Theorem~\ref{thm:main} and the definition of
$\Good_{\mathrm{low}}(\lambda)$.
\end{proof}

\section{Thinness principles that are refuted}\label{sec:consequences}

\begin{corollary}[No definable thin complete section]\label{cor:thin}
At an ultraexacting $\lambda$, no $T\in\OD_{V_\lambda}$ is a thin complete
section of $Q_\lambda$. In particular this excludes transversals,
finite-valued complete sections, countable-valued complete sections, and
sections whose fiber sizes vary below $\lambda$ without a common bound.
\end{corollary}

\begin{proof}
Pick $q\in\Good_{\mathrm{low}}(\lambda)$. Completeness makes $T\cap q$
nonempty, and goodness makes its cardinality $\lambda$, contradicting
thinness.
\end{proof}

\begin{corollary}[A common saturation set]\label{cor:common}
Every $\OD_{V_\lambda}$ complete section $T$ satisfies
\[
      \forall q\in\Good_{\mathrm{low}}(\lambda)\quad |T\cap q|=\lambda.
\]
The same set $\Good_{\mathrm{low}}(\lambda)$ works for all such $T$ and has
cardinality at least $\lambda$.
\end{corollary}

\begin{proof}
This is the definition of goodness together with completeness and
Theorem~\ref{thm:main}.
\end{proof}

\begin{corollary}[No local definable thin coverage]\label{cor:local-thin}
The following apparently weaker requirement fails:
\[
 \forall q\in Q_\lambda\ \exists T\in\OD_{V_\lambda}\quad
       (T\subseteq D_\lambda\ \wedge\ 0<|T\cap q|<\lambda).
\]
The set $T$ is allowed to depend on $q$ and need not be a complete section.
\end{corollary}

\begin{proof}
Every member of $\Good_{\mathrm{low}}(\lambda)$ witnesses failure.
\end{proof}

\begin{corollary}[Ordinal-valued maps have large fibers]\label{cor:coloring}
If $f:D_\lambda\to\Ord$ is a set function in $\OD_{V_\lambda}$, then for every
$q\in\Good_{\mathrm{low}}(\lambda)$ and every ordinal $\xi$,
\[
        |q\cap f^{-1}\{\xi\}|\in\{0,\lambda\}.
\]
Thus there is no $\OD_{V_\lambda}$ ordinal-valued map whose nonempty fibers
inside every equivalence class all have cardinality below $\lambda$.
\end{corollary}

\begin{proof}
The set $T_\xi=f^{-1}\{\xi\}$ belongs to $\OD_{V_\lambda}$: the additional parameter
$\xi$ is ordinal and is allowed. Apply the definition of $\Good_{\mathrm{low}}(\lambda)$.
For the last claim, any $q$ has at least one value of $f$, so at least one of
these fibers is nonempty.
\end{proof}

\subsection{Sharp cardinal and parameter cutoffs}

\begin{proposition}[The cardinal cutoff is sharp]\label{prop:cutoff}
The bound $<\lambda$ in Corollary~\ref{cor:thin} cannot be replaced by
$\le\lambda$: $T=D_\lambda$ is an ordinal-definable complete section and
$|T\cap q|=\lambda$ for every $q\in Q_\lambda$.
\end{proposition}

\begin{proof}
The set $D_\lambda$ is definable from the ordinal $\lambda$, and the size
assertion is Lemma~\ref{lem:class-size}.
\end{proof}

\begin{proposition}[The uniform parameter-rank cutoff is sharp]\label{prop:parameter-cutoff}
No class $q\in Q_\lambda$ has only empty or $\lambda$-sized intersections
with all $\OD_{V_{\lambda+1}}$ subsets of $D_\lambda$. Equivalently,
\[
                       \bigcap_{p\in V_{\lambda+1}}\Good_p(\lambda)=\varnothing.
\]
Thus $V_\lambda$ in the simultaneous theorem cannot be replaced by
$V_{\lambda+1}$.
\end{proposition}

\begin{proof}
For any $q\in Q_\lambda$, choose $a\in q$. The subset $a\subseteq\lambda$
belongs to $V_{\lambda+1}$. Taking the allowed parameter $p=a$, the set
$T=\{a\}$ belongs to $\OD_p$ and satisfies $|T\cap q|=1<\lambda$.
Hence $q\notin\Good_p(\lambda)$.
\end{proof}

\section{Small families without a uniform fiber bound}\label{sec:families}

The preceding results concern a single definable set of representatives. We
now allow the sections themselves to be nondefinable. Only the family
containing them is required to be ordinal-definable from low-rank parameters.
The new difficulty is that a union of fewer than $\lambda$ many sets of size
below $\lambda$ can have size $\lambda$, since $\cf(\lambda)=\omega$.
The embedding itself supplies the bound that ambient cardinal arithmetic
alone does not supply.

\begin{theorem}[Local obstruction to a small family of thin sections]\label{thm:local-families}
Let $j:X\to V_\zeta$ be an ultraexacting witness at $\lambda$, with
$\zeta>\lambda+\omega$. There is no fixed set $\mathcal F\in X$ such that
\[
 \varnothing\ne\mathcal F,\qquad |\mathcal F|<\lambda,
 \qquad \text{every }T\in\mathcal F\text{ is a thin complete section of }Q_\lambda.
\]
No assumption is made that the members of $\mathcal F$ belong individually to
$X$, are fixed, or have their fibers bounded by one common cardinal.
\end{theorem}

\begin{proof}
Let $\kappa=\crit(j)$ and let $q=q_0$ be the fixed critical-sequence class
from Lemma~\ref{lem:shifted}. Form the set of actual cardinal numbers
\[
       A=\{|T\cap q|:T\in\mathcal F\}.
\]
Thinness gives $A\subseteq\lambda$, and $|A|\le|\mathcal F|<\lambda$.
All the sets and possible bijections needed to compute these cardinalities
have rank below $\zeta$. Consequently $A$ is uniquely definable in $V_\zeta$
from $\mathcal F,q$, so $A\in X$. Since both these parameters are fixed,
$j(A)=A$. Lemma~\ref{lem:fixed-ordinals} gives
\begin{equation}\label{eq:fiber-bound}
                       A\subseteq\kappa.
\end{equation}
Thus a uniform bound has been obtained \emph{at this class}: every
$T\in\mathcal F$ has $|T\cap q|<\kappa$.

Now put
\[
              B=q\cap\bigcup\mathcal F
                =\bigcup_{T\in\mathcal F}(T\cap q).
\]
This is another fixed member of $X$. It is nonempty, because $\mathcal F$
is nonempty and every one of its sections meets $q$. It is a family of
cofinal $\omega$-subsets of $\lambda$. In the ambient universe, writing
$\mu=|\mathcal F|<\lambda$, Choice and \eqref{eq:fiber-bound} give
\[
                  |B|\le\mu\cdot\kappa<\lambda.
\]
The last inequality uses only that $\mu,\kappa<\lambda$ and that $\kappa$
is infinite: their product is at most $\max\{\mu,\kappa\}$ when $\mu$ is
infinite, and at most $\kappa$ otherwise. We have contradicted
Lemma~\ref{lem:small-family}.
\end{proof}

\begin{theorem}[No definable small family of thin sections]\label{thm:families}
If $\lambda$ is ultraexacting, there is no nonempty family
$\mathcal F\in\OD_{V_\lambda}$ of cardinality below $\lambda$ all of whose
members are thin complete sections of $Q_\lambda$.
\end{theorem}

\begin{proof}
Suppose such a family exists, with $\mathcal F\in\OD_p$ for some
$p\in V_\lambda$. Every candidate family is a subset of $\Pow(D_\lambda)$,
so all candidates lie in $V_{\lambda+12}$. Stabilize $\OD_p$ and its canonical
order on this rank using Lemma~\ref{lem:stabilization} with $P=\{p\}$, and choose an
ultraexacting witness at the resulting sufficiently large height, with
critical point above $\rank(p)$.

The property ``nonempty, of size below $\lambda$, and consisting of thin
complete sections'' is correct at this height: it quantifies over the full
$D_\lambda$ and $Q_\lambda$, and its cardinality comparisons use bijections
of bounded rank. The canonical least $\OD_p$ family having that property is
therefore a fixed member of $X$, by the proof of
Lemma~\ref{lem:canonical} with parameters $p,\lambda$. This contradicts
Theorem~\ref{thm:local-families}.
\end{proof}

\begin{corollary}[A stronger inconsistent theory]\label{cor:family-inconsistency}
The following theory is inconsistent:
\[
\begin{gathered}
 \ZFC+\text{``$\lambda$ is ultraexacting''}\\
 {}+\text{``there is a nonempty $\OD_{V_\lambda}$ family of fewer than $\lambda$}\\
 \text{thin complete sections of $Q_\lambda$''.}
\end{gathered}
\]
In particular, a countable definable family of possibly nondefinable
transversals is impossible.
\end{corollary}

\begin{proof}
The theorem excludes precisely the displayed conjunction. A transversal is
thin, and a countable family has size below the uncountable $\lambda$.
\end{proof}

\begin{remark}[Which earlier question this does not settle]\label{rem:selector-distinction}
A thin complete section is a subset of the \emph{representative space}
$D_\lambda$. A selector on finite subsets of $Q_\lambda$ chooses one
\emph{equivalence class} from each finite collection of classes. These are
different types of objects. The synthesis's question about a countable or
$<\lambda$-sized definable family of such selectors is not resolved by
Theorem~\ref{thm:families}. No reduction from that problem to the present one
has been proved.
\end{remark}

\section{A preserved thin predicate obstructs an Icarus embedding}\label{sec:icarus}

The local argument does not require the whole exacting framework. If an inner
model already contains every subset of $V_\lambda$, the critical sequence is
available there without an ultraexacting domain hypothesis. This gives a
precise obstruction to enriching a rank-into-rank inner model by a thin
complete section that the embedding must preserve.

\needspace{13\baselineskip}
\begin{theorem}[Thin-predicate obstruction]\label{thm:icarus}
Work in an ambient universe satisfying $\ZFC$. Let $\lambda$ be an uncountable cardinal, and let $N$ be a transitive inner
model of $\ZF$ such that $V_{\lambda+1}\subseteq N$. Suppose there is an
elementary embedding $j:N\to N$ with critical point $\kappa<\lambda$ whose
critical sequence is a set and satisfies
\[
                \sup_{n<\omega}j^n(\kappa)=\lambda.
\]
Then $j$ cannot preserve any thin complete section $T\in N$ of $Q_\lambda$:
\[
                     \Thin(T,\lambda)\quad\Longrightarrow\quad j(T)\ne T.
\]
Thinness and its cardinal bounds are measured in the ambient universe.
No form of Choice in $N$ is assumed.
\end{theorem}

\begin{proof}
Write $\kappa_n=j^n(\kappa)$. The sequence
$c=\langle\kappa_n:n<\omega\rangle$ is a subset of $V_\lambda$, hence belongs
to $V_{\lambda+1}$ and to $N$. Its range
$a=\{\kappa_n:n<\omega\}$ belongs to $N$ as well. Elementarity and the fact
that $j$ fixes $\omega$ give
\[
 j(c)(n)=\kappa_{n+1},\qquad j(a)=a\setminus\{\kappa\},
 \qquad j(\lambda)=\lambda.
\]
The last equality follows by applying $j$ to $\lambda=\sup\ran(c)$ inside
$N$.

The model $N$ computes $D_\lambda$ correctly. Indeed, it contains all subsets
of $\lambda$; order type $\omega$ and cofinality in this definition are
absolute for these wellordered subsets. Their increasing enumerations are
uniquely definable and belong to $N$. Thus $N$ also computes the full class
$q=[a]$ correctly, and $q\in N$. Finite modification gives $j(q)=q$.

Suppose $T\in N$ is a thin complete section and $j(T)=T$. Then
\[
                 B=T\cap q,\qquad A=\bigcup B
\]
are fixed sets in $N$. Ambient thinness makes $B$ a nonempty family of fewer
than $\lambda$ countable sets. Ambient Choice therefore gives $|A|<\lambda$,
while $A$ is cofinal in $\lambda$.

Let $\tau=\ot(A)$, computed in either transitive universe; this ordinal is
absolute. Since $\lambda$ is a cardinal and $|A|<\lambda$, we have
$\tau<\lambda$. The increasing enumeration $h:\tau\to A$ belongs to $N$,
and $j(A)=A$ implies $j(\tau)=\tau$ and $j(h)=h$. Every ordinal
$\xi\in[\kappa,\lambda)$ is moved: if
$\kappa_n\le\xi<\kappa_{n+1}$, then
$j(\xi)\ge\kappa_{n+1}>\xi$. Hence $\tau<\kappa$. For $\xi<\tau$,
\[
                    j(h(\xi))=j(h)(j(\xi))=h(\xi),
\]
so every element of $A$ lies below $\kappa$. This contradicts its cofinality.
\end{proof}

\begin{corollary}[An explicit forbidden enrichment]\label{cor:icarus-enrichment}
Let $T$ be any actual thin complete section of $Q_\lambda$. There is no
elementary self-embedding of the structure
\[
                       (L(V_{\lambda+1},T),\in,T)
\]
with critical point below $\lambda$ and critical sequence cofinal in $\lambda$.
Here the final $T$ is a named set, so its preservation is part of the
structure's elementarity requirement.
\end{corollary}

\begin{proof}
Take $N=L(V_{\lambda+1},T)$, understood as the constructible inner model
containing $V_{\lambda+1}$ and $T$ as sets. Its transitivity gives
$V_{\lambda+1}\subseteq N$, and naming $T$ requires $j(T)=T$.
Theorem~\ref{thm:icarus} applies.
\end{proof}

There are thin complete sections in the ambient universe by Choice: choose
one representative from each member of the set $Q_\lambda$. The corollary
therefore excludes genuine enrichments, not an empty collection of possible
parameters. On the other hand, the preservation requirement is essential to
the argument. An embedding of the underlying inner model that moves $T$ is
not refuted. Nor does the argument refute $\Izero$, whose model need not
contain a thin section that is fixed by its embedding.

\section{A bounded-rank separation theorem from ordinary exactingness}\label{sec:separation}

The saturation theorem used ultraexactingness to put the critical sequence
inside the domain. A different consequence needs only exactingness. It
clarifies why a small definable family of objects of rank $\lambda+1$ cannot
hide its distinctions at ranks cofinal in $\lambda$.

\begin{lemma}[The ordinary exacting cofinality barrier]\label{lem:ordinary-barrier}
If $\lambda$ is exacting and $S\in\OD_{V_\lambda}$ is a subset of $\lambda$
with $|S|<\lambda$, then $\sup S<\lambda$.
\end{lemma}

\begin{proof}
Suppose instead that such a set is cofinal, and choose $p\in V_\lambda$
such that some cofinal $S\in\OD_p$ has size below $\lambda$. Stabilize the
canonical $\OD_p$ order on a rank containing all subsets of $\lambda$, and
choose an ordinary exacting witness with critical point above $\rank(p)$.
At this height the property ``a cofinal subset of $\lambda$ of cardinality
below $\lambda$'' is correct. Its canonical least $\OD_p$ witness is fixed
and belongs to $X$. Lemma~\ref{lem:fixed-ordinals} is valid for ordinary
exacting witnesses and supplies the contradiction.
\end{proof}

This is the familiar cofinality barrier, not a new large-cardinal comparison.
Equivalently, $\lambda$ is regular in $\HOD_{V_\lambda}$, as proved in
\cite[Theorem 2.10]{abl}. For the direction needed here, a cofinal sequence in
that inner model would have an $\OD_{V_\lambda}$ range of size below
$\lambda$. We included the argument to make the separation proof independent
of an unexplained appeal to inner-model regularity.

\begin{theorem}[One bounded rank separates a small definable family]\label{thm:separation}
Suppose $\lambda$ is exacting,
\[
       A\in\OD_{V_\lambda},\qquad A\subseteq\Pow(V_\lambda)=V_{\lambda+1},
       \qquad |A|<\lambda.
\]
Then some $\beta<\lambda$ makes the restriction map
\[
                   r_\beta:A\longrightarrow\Pow(V_\beta),
                   \qquad r_\beta(x)=x\cap V_\beta
\]
injective. Moreover,
\[
       A\in\HOD_{V_\lambda},
       \qquad \text{and $A$ has an enumeration of length $|A|$ in }
                       \HOD_{V_\lambda}.
\]
Here $|A|$ is its actual ambient cardinality, not a potentially larger
inner-model cardinality.
\end{theorem}

\begin{proof}
For distinct $x,y\in A$, let
\[
 d(x,y)=\min\{\alpha<\lambda:x\cap V_\alpha\ne y\cap V_\alpha\}.
\]
This minimum exists because $x,y\subseteq V_\lambda$: a member of their
symmetric difference lies in some bounded rank. Put
\[
                    S=\{d(x,y):x,y\in A,\ x\ne y\}.
\]
The definition of $S$ uses $A$ and the ordinal $\lambda$, so
$S\in\OD_{V_\lambda}$. Also $|S|\le|A|^2<\lambda$.
Lemma~\ref{lem:ordinary-barrier} makes $S$ bounded in $\lambda$.
Choose $\beta<\lambda$ above all members of $S$; when $S$ is empty, any
$\beta<\lambda$ suffices. If $x\ne y$, their restrictions already differ at
$d(x,y)$ and still differ at $\beta$. This proves injectivity.

Write $H=\HOD_{V_\lambda}$. For $x\in A$, its restriction
$r=x\cap V_\beta$ belongs to $V_{\beta+1}\subseteq V_\lambda$. The object $x$
is the unique member of $A$ with this restriction. A definition of $A$ from
low-rank parameters, together with the new low-rank parameter $r$ and the
ordinal $\beta$, therefore defines $x$. Thus $x\in\OD_{V_\lambda}$.
Every element of $x$ belongs to $V_\lambda$, and all of its transitive closure
does too. These sets are hereditarily definable from allowed low-rank
parameters. Hence $x\in H$. Since $A$ itself is ordinal-definable from such
parameters and all its members are hereditary in this sense, $A\in H$.

Choose in $V$ a wellordering $w$ of $\Pow(V_\beta)$. Its rank is below
$\lambda$, so $w\in V_\lambda\subseteq H$. Pull it back along the injective
map $r_\beta$ to obtain a wellordering $\triangleleft$ of $A$ in $H$.
The order type $\tau$ is the actual order type, by transitivity, and has
cardinality $|A|<\lambda$. Since $\lambda$ is a cardinal, $\tau<\lambda$.
The increasing enumeration $h:\tau\to A$ is computed correctly in $H$.
Finally let $\mu=|A|$. An ambient bijection $b:\mu\to\tau$ has rank below
$\lambda$, hence belongs to $V_\lambda\subseteq H$. The composition
$h\circ b$ is an enumeration of $A$ of length $\mu$ in $H$.
The same argument handles the empty family with the empty enumeration.
\end{proof}

The final proof does not assume that $H$ satisfies global Choice. It constructs
the required wellordering from a particular bounded-rank wellordering that
is already an element of $V_\lambda$. The rank restriction on the members of
$A$ is also substantive: a selector defined on arbitrary equivalence classes
has a graph of higher rank, and the above first-difference map need not take
values below $\lambda$. This is another reason the unresolved selector problem
does not follow from the separation theorem.

\section{An \texorpdfstring{$\Izero$}{I0}-equiconsistent uniform-saturation package}\label{sec:consistency}

The supplied reports already identify the consistency of ultraexactingness
with a completely ordinal-definable rank below $\lambda$. The new theorem
adds simultaneous low-parameter saturation to that model. Unlike the
positive coding hypothesis, this saturation follows from ultraexactingness
alone. Coding the lower rank into $\HOD$ then supplies a useful contrasting
region where definable choice remains available.

\begin{lemma}[Collapse of the parameter hierarchy]\label{lem:parameter-collapse}
If $V_\lambda\subseteq\HOD$, then
\[
             \OD_{V_\lambda}=\OD_p=\OD
                         \qquad\text{for every }p\in V_\lambda.
\]
If $\lambda$ is also ultraexacting, then
\[
        \Good_{\mathrm{low}}(\lambda)=\Good_p(\lambda)=\Good_{\varnothing}(\lambda),
        \qquad |\Good_{\mathrm{low}}(\lambda)|\ge\lambda.
\]
\end{lemma}

\begin{proof}
An allowed parameter $p$ is hereditarily ordinal-definable and therefore is
itself ordinal-definable. Substitute a definition of $p$ from its finite
ordinal tuple into any definition using $p$ and ordinal parameters. The
result uses only finitely many ordinal parameters. The same substitution
works for a finite list of elements of $V_\lambda$. This proves containment
in $\OD$; reverse containment is immediate by ignoring the extra parameters.
The equalities of the three sets of classes follow from their definitions.
The cardinal bound follows from Theorem~\ref{thm:main}.
\end{proof}

\subsection{A positive control below the cofinality boundary}

Put
\[
 B_\lambda=\{a\subseteq\lambda: |a|=\aleph_0\text{ and }\sup a<\lambda\},
 \qquad Q^{\mathrm{bd}}_\lambda=B_\lambda/{\eqfin}.
\]
The members of $B_\lambda$ are all bounded, in contrast to $D_\lambda$.
They need not have order type $\omega$: the displayed definition uses
countable infinitude, which is sufficient for the control example.

\begin{proposition}[The bounded quotient still has definable choice]\label{prop:bounded-positive}
If $\lambda$ is an uncountable limit cardinal and $V_\lambda\subseteq\HOD$,
then $Q^{\mathrm{bd}}_\lambda$ has an ordinal-definable wellordering and an
ordinal-definable transversal. These statements concern the full ambient
bounded quotient, not merely a quotient computed using fewer representatives
in $\HOD$.
\end{proposition}

\begin{proof}
Every $a\in B_\lambda$ has rank below $\lambda$, so $a\in\HOD$. Its
countability has a witness of rank below $\lambda$: an enumeration of a
bounded subset of $\lambda$ by $\omega$ has bounded rank. Such a witness
also belongs to $\HOD$. Thus $\HOD$ recognizes exactly the same members of
$B_\lambda$. Moreover $B_\lambda$ is ordinal-definable from $\lambda$ and
all its elements are hereditary in $\HOD$, so $B_\lambda\in\HOD$.

For $a\in B_\lambda$, let $[a]_{\mathrm{bd}}$ be its finite-modification
class within $B_\lambda$. Every member of this class is bounded: finitely
many additions to a bounded set remain bounded below the limit cardinal
$\lambda$. The full class is definable from $a,\lambda$ and consists of
members of $\HOD$. Since $a$ is ordinal-definable, so is the class; it
therefore belongs to $\HOD$. The full quotient $Q^{\mathrm{bd}}_\lambda$
is ordinal-definable from $\lambda$ and all its members belong to $\HOD$,
so it too belongs to $\HOD$.

Restrict the canonical definable wellorder of $\HOD$ to the set
$Q^{\mathrm{bd}}_\lambda$. This gives an ordinal-definable wellordering.
For each class, take its least representative in the canonical wellorder
of $\HOD$. Each class is nonempty and consists of $\HOD$ sets, so the least
representative exists. The set of all these representatives is definable
from the ordinal $\lambda$ and meets each class exactly once. It is the
required transversal.
\end{proof}

This control separates two issues. The classes of the bounded quotient also
have size $\lambda$, by the finite-addition argument of
Lemma~\ref{lem:class-size}; their large size alone does not cause the
obstruction. The obstruction uses the cofinality of \emph{each representative}
and the resulting cofinality of a fixed small union.

\subsection{The formal comparison}

Let $\Tbase$ be the theory
\[
           \ZFC+\exists\lambda\,
                 \bigl(\UEx(\lambda)\wedge V_\lambda\subseteq\HOD\bigr).
\]
Let $\Tsat$ assert that there is a cardinal $\lambda$ with all the following
properties, in addition to $\ZFC$:
\begin{enumerate}[label=\textup{(\roman*)}]
\item $\lambda$ is ultraexacting and $V_\lambda\subseteq\HOD$;
\item $|\Good_{\mathrm{low}}(\lambda)|\ge\lambda$;
\item there is no nonempty $\OD_{V_\lambda}$ family of fewer than $\lambda$
      thin complete sections of $Q_\lambda$;
\item the full bounded quotient $Q^{\mathrm{bd}}_\lambda$ has an
      ordinal-definable wellordering and an ordinal-definable transversal.
\end{enumerate}
All these assertions are first-order expressible. In particular, local
satisfaction codes express ordinal definability, and set embeddings express
ultraexactingness. The usual $\Izero$ theory is understood in the formulation
used in Input~\ref{in:cons}.

\begin{theorem}[Equiconsistency with simultaneous low-parameter saturation]\label{thm:equiconsistency}
The theories $\Tsat$ and $\Tbase$ are equivalent over $\ZFC$, and
\begin{equation}\label{eq:equiconsistency}
           \Con(\Tsat)\quad\Longleftrightarrow\quad\Con(\ZFC+\Izero).
\end{equation}
This is a refinement of the properties of a known consistency model,
not a claim to lower or newly determine the consistency strength of
ultraexactingness.
\end{theorem}

\begin{proof}
First, any model of $\Tsat$ satisfies $\Tbase$ by clause (i). Conversely,
let $\lambda$ witness $\Tbase$. Theorem~\ref{thm:main} gives
clause (ii); Theorem~\ref{thm:families} gives (iii); and
Proposition~\ref{prop:bounded-positive} gives (iv). The same $\lambda$
witnesses all clauses. Thus the two theories are equivalent.

For the forward consistency implication in \eqref{eq:equiconsistency},
$\Tsat$ implies the existence of an ultraexacting cardinal. The
ultraexacting-to-$\Izero$ direction of Input~\ref{in:cons} then yields
$\Con(\ZFC+\Izero)$.

For the reverse implication, the other direction of Input~\ref{in:cons}
first gives the consistency of $\ZFC$ with an ultraexacting cardinal.
Its preservation-and-coding clause supplies the relative consistency of
an ultraexacting $\lambda$ with $V_\lambda\subseteq\HOD$, namely
$\Con(\Tbase)$. The proved implication $\Tbase\vdash\Tsat$ yields
$\Con(\Tsat)$. These are the standard syntactic relative-consistency
implications obtained from the cited comparison and forcing theorem; no
transitive-model assumption is added to the hypothesis $\Con(\ZFC+\Izero)$.
\end{proof}

\begin{corollary}[Consistency of genuinely uniform classes]\label{cor:all-low-consistent}
Assuming $\Con(\ZFC+\Izero)$, it is consistent that an ultraexacting
$\lambda$ has at least $\lambda$ classes $q$ such that \emph{every}
$\OD_{V_\lambda}$ subset of $D_\lambda$ meets $q$ in either zero or $\lambda$
representatives. At the same time the bounded quotient admits a definable
transversal, and the cofinal quotient admits transversals by Choice but
none that is $\OD_{V_\lambda}$.
\end{corollary}

\begin{proof}
Take a model of $\Tsat$. Clause (ii) is exactly the first assertion, and
clause (iv) gives the bounded transversal. Ambient Choice selects a
representative from every member of the set $Q_\lambda$.
Corollary~\ref{cor:thin} excludes a low-parameter ordinal definition of
any such transversal.
\end{proof}

\needspace{8\baselineskip}
\section{Proof audit and limits of the result}\label{sec:audit}

\subsection{The dependency graph}

The essential proof has the following short chain:
\[
\begin{gathered}
 \text{local Kunen + exacting witness}
 \ \Longrightarrow\ \text{no fixed small cofinal subset of }\lambda,\\
 \text{no fixed small cofinal subset}
 \ \Longrightarrow\ \text{no fixed small family in }D_\lambda,\\
 \text{ultraexactingness}
 \ \Longrightarrow\ \text{fixed lower-rank code and many fixed classes},\\
 \text{bounded-rank definition codes + fixed classes}
 \ \Longrightarrow\ \text{simultaneous full-sized traces}.
\end{gathered}
\]
The family theorem adds only the fixed set of fiber cardinalities. The
separation theorem uses only the ordinary exacting cofinality barrier.
The Icarus obstruction is a separate application to an inner model with the
critical sequence already present. The equiconsistency refinement imports
Input~\ref{in:cons}, after all the new combinatorial implications have been
proved in $\ZFC$.

\subsection{Checks at the main potential failure points}

\textbf{Ordinal parameters.}
We never infer that an arbitrary $\OD_p$ set is fixed merely because $p$ is
fixed. Its ordinal parameters might be moved. Instead, bounded-rank
stabilization is arranged before the embedding is chosen, and the least
counterexample is defined from parameters known to be fixed.

\textbf{A nontransitive domain.}
Membership $B\in X$ does not imply $B\subseteq X$. The small-family proofs
do not make that inference. They use definable union and, for a set of
ordinals, the uniquely definable increasing enumeration. Pointwise image
arguments are used only when an actual enumeration by fixed natural numbers
belongs to $X$.

\textbf{Singular cardinal arithmetic.}
For a small family of countable sets, the estimate is $\mu\cdot\aleph_0$.
For a small family of thin sections, the bound $\mu\cdot\kappa$ is justified
only after the embedding has forced all relevant fiber cardinals below
$\kappa$. Replacing this step by ``a union of $<\lambda$ small sets is small''
would be invalid.

\textbf{Correctness chosen in the right order.}
The height $\zeta$ is first made correct for all relevant bounded-rank
$\OD_p$ objects, uniformly over $p\in V_{\lambda+1}$. The embedding, its
fixed code $b$, and the classes $q_t$ are selected afterwards. Neither the
critical sequence nor the coding parameter is used prematurely to choose
the correctness height.

\textbf{External versus internal cardinalities.}
All smallness estimates use ambient $\ZFC$. In the Icarus argument, the
inner model need not see a bijection witnessing the smallness of $B$.
It needs only the order type of $\bigcup B$, which is absolute and lies
below $\lambda$ by the ambient estimate.

\textbf{Equiconsistency versus implication in the same universe.}
The proof does not assert $\UEx(\lambda)\Rightarrow\Izero(\lambda)$.
It uses an equiconsistency theorem, followed in the appropriate direction by
a preservation theorem, and explicitly proves that the new clauses hold in
the resulting model.

\subsection{What remains open in this continuation}

The simultaneous conclusion over $V_\lambda$ is settled by
Theorem~\ref{thm:main}; it does not require $V_\lambda\subseteq\HOD$.
Proposition~\ref{prop:parameter-cutoff} settles the next uniform rank in the
opposite direction. The remaining questions are genuinely different.

\begin{question}
Can the lower bound $|\Good_{\mathrm{low}}(\lambda)|\ge\lambda$ be strengthened, for
example to $2^\lambda$, under the same hypothesis? The proof produces
arbitrarily large critical-point-sized families below $\lambda$; it does not
establish that its bound is optimal.
\end{question}

\begin{question}
Can ordinary exactingness replace ultraexactingness in the thin-section
obstruction? The local obstruction already needs only exactingness, but the
construction of a fixed critical-sequence class uses
$j\restr V_\lambda\in X$. An alternative way to obtain a fixed suitable
class would be needed.
\end{question}

The supplied synthesis's small-family problem for selectors on finite
subsets of $Q_\lambda$ remains distinct and unresolved here, as does its
cover-exacting/strongly-compact interaction problem. No conclusion about
those questions is inferred from terminology shared with the present
thin-section results.

\subsection{Novelty and verification status}

The contribution relative to the supplied documents is the replacement of
finite fibers by arbitrary fibers of size $<\lambda$, the simultaneous
$\lambda$-many-class conclusion over all low-rank parameters, the nonuniform small-family extension, the
explicit preserved-predicate obstruction, and the stated rank-separation
lemma. The low-parameter saturation theorem, which needs only
ultraexactingness, is then added to the
existing $\Izero$-equiconsistent coding model.

The primary-source checks support the external inputs and distinguish them
from the new deductions. They do not certify worldwide priority: related
consequences may be known to specialists or appear under different
formulations. The proofs have been checked internally for the issues
listed above, but have not been formalized in a proof assistant or
independently refereed. These qualifications apply to the new results, not
to a claimed proof of an unconditional inconsistency, which this report does
not present.

\tcbset{breakable=false}
\begin{assessment}{Research conclusion}
The finite-transversal boundary can be strengthened to a cardinal-sized
one: at an ultraexacting $\lambda$, a definable complete section must have
$\lambda$ representatives in many classes, and the classes can be chosen
simultaneously for all definitions with arbitrary parameters from
$V_\lambda$. The parameter bound is sharp at the next rank.
Even a definable small family of nondefinable thin sections is impossible.

At the same consistency strength as $\Izero$, this uniform saturation
coexists with $V_\lambda\subseteq\HOD$, while definable choice survives
on the bounded quotient. This yields both a conditional
inconsistency theorem and a positive relative-consistency refinement, with
no claim that ultraexactingness itself is inconsistent.
\end{assessment}

\appendix
\section{Rank and definability bookkeeping}\label{app:ranks}

This appendix records bounds used implicitly in the elementarity arguments.
They are intentionally generous; optimizing the finite rank offset would
not change any conclusion.

For an infinite limit cardinal $\lambda$, each ordinal below $\lambda$ and
every finite tuple of such ordinals has rank below $\lambda$. A subset of
$\lambda$ has rank at most $\lambda$, hence belongs to $V_{\lambda+1}$.
It follows that
\[
\begin{array}{c|c}
\text{Object} & \text{A rank containing it}\\ \hline
 a\subseteq\lambda & V_{\lambda+1}\\
 D_\lambda,\ q\in Q_\lambda,\ T\subseteq D_\lambda & V_{\lambda+2}\\
 Q_\lambda,\ \mathcal F\subseteq\Pow(D_\lambda) & V_{\lambda+3}
\end{array}
\]
A function whose domain is an ordinal below $\lambda$ and whose values
are members of any one of these displayed sets has graph of rank at most
$\lambda$ plus a fixed finite number. In particular, all possible bijections
witnessing cardinalities $<\lambda$, all increasing enumerations of subsets
of $\lambda$, and the maps used to list fiber cardinalities occur below
$V_{\lambda+12}$. A height $\zeta>\lambda+\omega$ therefore contains the full
collection of relevant witnesses. This proves the cardinality absoluteness
used in Sections~\ref{sec:od} and~\ref{sec:families}.

The relative hereditary class $\HOD_{V_\lambda}$ is the class of sets all of
whose transitive-closure members, including the set itself, belong to
$\OD_{V_\lambda}$. It is transitive, contains $V_\lambda$, and satisfies
$\ZF$. For example, its internal power set of $x$ is defined by separation
from $\Pow(x)$ using the definable hereditary class; its internal replacement
sets are obtained by ambient Replacement for the relativized formula.
The resulting sets are definable from the finitely many low-rank parameters
needed for $x$ and the formula's parameters, together with the ordinal
$\lambda$, and their elements remain hereditary. Finite lists of low-rank
parameters can be combined into one because $\lambda$ is a limit cardinal.
These observations justify the closure operations used in
Theorem~\ref{thm:separation}; they do not assert global Choice in that inner
model.

Finally, the canonical code order in Section~\ref{sec:od} can be made
completely explicit. At a fixed height $\theta$, order codes first by the
finite tuple length, then by the formula number, then lexicographically on
the ordinal tuple. For a fixed finite length this is a wellordering. The
ordinal sum of these orders over the natural numbers is still a set
wellordering. Placing the height $\theta$ first gives a definable class order
with a set of predecessors below every code. Thus ``the least code'' does
not appeal to an arbitrary global choice function.

\section{Build and source notes}\label{app:build}

The accompanying source is self-contained apart from standard TeX Live
packages. It uses the same \texttt{newpxtext}/\texttt{newpxmath} font family,
page geometry, heading treatment, callout boxes, and forest/olive/sage palette
as the supplied reports. The colors are \texttt{30392C}, \texttt{687144},
\texttt{65685F}, \texttt{CBD0BC}, and \texttt{F2F4EB}.
The source can be compiled with two or three runs of \texttt{pdflatex}; the
archive includes build instructions. No font files are included.

The bibliography distinguishes the two supplied manuscripts from primary
external sources. The references to the supplied synthesis concern its
finite-modification quotient and finite-valued transversal arguments; the
new proofs are written in full here and do not depend on its unpublished
proofs. The source versions below were checked on 18 September 2026.

\begin{thebibliography}{9}
\addcontentsline{toc}{section}{References}

\bibitem{abl}
Juan P. Aguilera, Joan Bagaria, and Philipp L\"ucke.
\emph{Large cardinals, structural reflection, and the HOD Conjecture}.
\arxiv{2411.11568}. Version consulted:
\url{https://arxiv.org/html/2411.11568v4}.
Especially Lemmas 2.3 and 3.2 and Theorem 2.10.

\bibitem{abgl}
Juan P. Aguilera, Joan Bagaria, Gabriel Goldberg, and Philipp L\"ucke.
\emph{Large cardinals beyond HOD}.
\arxiv{2509.10254}. HTML version consulted:
\url{https://arxiv.org/html/2509.10254v1}.
Especially Proposition 3.9 and Theorem 4.5.

\bibitem{kunen}
Kenneth Kunen.
Elementary embeddings and infinitary combinatorics.
\emph{Journal of Symbolic Logic} \textbf{36}(3) (1971), 407--413.
\doi{10.2307/2269948}.

\bibitem{cover}
Gabriel Goldberg, joint work with Douglas Blue.
\emph{Consistency beyond Kunen inconsistency}.
Lecture slides, 1 July 2026. Especially Theorem 1.1 and Remark 2.4.
\url{https://math.berkeley.edu/~goldberg/Slides/CoverExact.pdf}.

\bibitem{synthesis}
\emph{Large cardinals at the inconsistency frontier: a synthesis}.
User-supplied research manuscript, 18 September 2026.
File: \nolinkurl{Large_Cardinals_Synthesis.tex} in \texttt{Cardinals2.zip}.
Especially the finite-choice and finite-transversal development in
Sections 8--9, and the subsequent consistency discussion.

\bibitem{unified}
\emph{Large cardinals at the inconsistency frontier: a unified assessment}.
User-supplied research manuscript, 18 September 2026.
File: \nolinkurl{Large_Cardinals_Unified_Report.tex} in \texttt{Cardinals2.zip}.

\end{thebibliography}
\end{document}
